
\documentclass{article}
\usepackage{colortbl}
\usepackage{makecell}
\usepackage{multirow}
\usepackage{supertabular}

\begin{document}

\newcounter{utterance}

\centering \large Interaction Transcript for game `eqbench', experiment `eq\_bench', episode 17 with reflexion{-}llama{-}3.1{-}8b{-}instruct.
\vspace{24pt}

{ \footnotesize  \setcounter{utterance}{1}
\setlength{\tabcolsep}{0pt}
\begin{supertabular}{c@{$\;$}|p{.15\linewidth}@{}p{.15\linewidth}p{.15\linewidth}p{.15\linewidth}p{.15\linewidth}p{.15\linewidth}}
    \# & \multicolumn{2}{c}{Player} && \multicolumn{2}{c}{Game Master} \\
    \hline

    \theutterance \stepcounter{utterance}  
    & & & \multicolumn{4}{p{0.6\linewidth}}{
        \cellcolor[rgb]{0.9,0.9,0.9}{
            \makecell[{{p{\linewidth}}}]{
                \texttt{\tiny{[P1$\langle$GM]}}
                \texttt{You have attempted this same task before and did not succeed. These are the plans you wrote after those attempts.} \\
\texttt{Trial \#1: New plan: } \\
\\ 
\texttt{Upon re{-}examining the dialogue, I realized that I focused too much on the surface{-}level exchange between Grayson and Dr. Sinclair, without considering the underlying dynamics and motivations of the characters. To improve my analysis, I will pay closer attention to the power dynamics, emotional cues, and the context of the conversation. Specifically, I will:} \\
\\ 
\texttt{1. Identify the underlying goals and motivations of each character, and how they relate to the conversation.} \\
\texttt{2. Analyze the emotional cues and body language (if present) to better understand the emotional state of each character.} \\
\texttt{3. Consider the context of the conversation, including the setting, the relationship between the characters, and any external factors that may influence their emotions.} \\
\\ 
\texttt{By taking a more nuanced and multi{-}faceted approach, I hope to improve my ability to predict the emotional responses of the characters and provide a more accurate answer.} \\
\\ 
\\ 
\\ 
\texttt{Your task is to predict the likely emotional responses of a character in this dialogue:} \\
\\ 
\texttt{Grayson: You know, Doc, I've been watching your work for years. All that time and money spent on Mars, and what have you got to show for it? } \\
\texttt{Dr. Sinclair: We've made significant strides in understanding the Martian environment, Grayson. It's not about instant results. } \\
\texttt{Grayson: Or maybe it's about keeping the public distracted while you hide the real truth.} \\
\texttt{Dr. Sinclair: And what truth would that be, Mr. Hawk?} \\
\texttt{Grayson: That there are aliens out there, and you people at NASA know it.} \\
\texttt{Dr. Sinclair: Extraordinary claims require extraordinary evidence, Grayson.} \\
\texttt{Grayson: Oh, I have evidence. I've got testimonials, photos, documents...} \\
\texttt{Dr. Sinclair: Anecdotes and blurry photos aren't evidence, Grayson. They're stories. } \\
\texttt{Grayson: Stories that you're too afraid to believe.} \\
\texttt{Dr. Sinclair: No, stories that don't hold up under scientific scrutiny. } \\
\texttt{Grayson: You're just a puppet, Sinclair. They've got you so convinced of your 'facts' that you can't see the truth.} \\
\texttt{Dr. Sinclair: And you're a sensationalist, Hawk. You've built a career on paranoia and misinformation.} \\
\texttt{Grayson: At least I'm not lying to the public.} \\
\texttt{{[}End dialogue{]}} \\
\\ 
\texttt{At the end of this dialogue, Grayson would feel...} \\
\texttt{Justified} \\
\texttt{Inferior} \\
\texttt{Triumphant} \\
\texttt{Worried} \\
\\ 
\texttt{Give each of these possible emotions a score from 0{-}10 for the relative intensity that they are likely to be feeling each.} \\
\\ 
\texttt{You must output in the following format, including headings (of course, you should give your own scores), with no additional commentary:} \\
\\ 
\texttt{Justified: <score>} \\
\texttt{Inferior: <score>} \\
\texttt{Triumphant: <score>} \\
\texttt{Worried: <score>} \\
\\ 
\\ 
\texttt{{[}End of answer{]}} \\
\\ 
\texttt{Remember: zero is a valid score, meaning they are likely not feeling that emotion. You must score at least one emotion > 0.} \\
\\ 
\texttt{Your answer:} \\
            }
        }
    }
     \\ \\

    \theutterance \stepcounter{utterance}  
    & \multicolumn{4}{p{0.6\linewidth}}{
        \cellcolor[rgb]{0.561,0.737,0.561}{
            \makecell[{{p{\linewidth}}}]{
                \texttt{\tiny{[P1$\rangle$GM]}}
                \texttt{Justified: 8} \\
\texttt{Inferior: 0} \\
\texttt{Triumphant: 9} \\
\texttt{Worried: 0} \\
            }
        }
    }
    & & \\ \\

    \theutterance \stepcounter{utterance}  
    & & & \multicolumn{2}{p{0.3\linewidth}}{
        \cellcolor[rgb]{0.9,0.9,0.9}{
            \makecell[{{p{\linewidth}}}]{
                \texttt{\tiny{[GM$|$GM]}}
                \texttt{Justified: 8} \\
\texttt{Inferior: 0} \\
\texttt{Triumphant: 9} \\
\texttt{Worried: 0} \\
            }
        }
    }
    & & \\ \\

    \theutterance \stepcounter{utterance}  
    & & & \multicolumn{2}{p{0.3\linewidth}}{
        \cellcolor[rgb]{0.9,0.9,0.9}{
            \makecell[{{p{\linewidth}}}]{
                \texttt{\tiny{[GM$|$GM]}}
                \texttt{Justified: 8} \\
\texttt{Inferior: 2} \\
\texttt{Triumphant: 6} \\
\texttt{Worried: 1} \\
            }
        }
    }
    & & \\ \\

    \theutterance \stepcounter{utterance}  
    & & & \multicolumn{2}{p{0.3\linewidth}}{
        \cellcolor[rgb]{0.9,0.9,0.9}{
            \makecell[{{p{\linewidth}}}]{
                \texttt{\tiny{[GM$|$GM]}}
                \texttt{game\_result = LOSE} \\
            }
        }
    }
    & & \\ \\

    \theutterance \stepcounter{utterance}  
    & & \multicolumn{4}{p{0.6\linewidth}}{
        \cellcolor[rgb]{0.95,0.95,0.95}{
            \makecell[{{p{\linewidth}}}]{
                \texttt{\tiny{[P1$\leftrightarrow$P]}}
                \texttt{Post{-}episode Reflexion (not sent to GM):} \\
\texttt{New plan: } \\
\\ 
\texttt{Upon re{-}examining the dialogue, I realized that I focused too much on the surface{-}level exchange between Grayson and Dr. Sinclair, without considering the underlying dynamics and motivations of the characters. To improve my analysis, I will pay closer attention to the power dynamics, emotional cues, and the context of the conversation. Specifically, I will:} \\
\\ 
\texttt{1. Identify the underlying goals and motivations of each character, and how they relate to the conversation. For example, Grayson's goal is to convince Dr. Sinclair of the existence of aliens, while Dr. Sinclair's goal is to maintain a scientific and rational approach.} \\
\texttt{2. Analyze the emotional cues and body language (if present) to better understand the emotional state of each character. For example, Grayson's tone and language suggest a sense of conviction and urgency, while Dr. Sinclair's tone and language suggest a sense of calm and skepticism.} \\
\texttt{3. Consider the context of the conversation, including the setting, the relationship between the characters, and any external factors that may influence their emotions. For example, the conversation takes place in a public setting, and Grayson's accusations may be motivated by a desire to gain attention and publicity.} \\
\\ 
\texttt{To apply this new plan, I will re{-}read the dialogue and look for evidence of these underlying dynamics and motivations. I will also consider alternative perspectives and emotional responses, and be prepared to adjust my answer accordingly.} \\
            }
        }
    }
    & \\ \\

\end{supertabular}
}

\end{document}
